\documentclass[11pt,a4paper]{article}
\usepackage[margin=2.5cm]{geometry}
\usepackage{amsmath,amssymb}
\usepackage{booktabs}
\usepackage{hyperref}
\usepackage{xcolor}
\newcommand{\code}[1]{\texttt{#1}}
\newcommand{\vect}[1]{\boldsymbol{#1}}

\title{Visualising GW250114 with SEOBNRv5PHM:\\how the posterior-cloud video is made}
\author{}
\date{October 2026}

\begin{document}
\maketitle

\begin{abstract}
We describe the pipeline that turns 200 posterior samples of the SEOBNRv5PHM
parameter-estimation run on GW250114 into a 3D animation of the binary's inspiral, merger
and ringdown, as seen from Earth, synchronised with the whitened LIGO Livingston strain.
Each sample's orbit, spins and remnant are computed from the model itself, called with the
same settings as the likelihood. We list what is physically faithful and what is
illustrative.
\end{abstract}

\section{Overview}

The video (\code{renders/cloud200.mp4}, $1920\times1080$, 30\,fps, 66\,s) has two parts:
a 3D view ($1920\times810$) rendered in Blender, and a panel ($1920\times270$) showing the
whitened L1 strain with a time cursor. The signal is played back 100 times slower than
real time, over the detector-time window $[-0.29, 0.37]$\,s around
$t_0 = 1420878141.0$ (merger at $\approx t_0 + 0.217$\,s), giving 1981 frames.

The pipeline, run from \code{gw250114/}, is:
\begin{enumerate}
  \item \code{validate\_dynamics.py}: consistency checks on the dynamics (Sec.~\ref{sec:checks});
  \item \code{ppd.py}: whitened posterior-predictive band for the panel (Sec.~\ref{sec:panel});
  \item \code{export\_scene.py --n-samples 200}: per-frame scene data for each sample
        (Secs.~\ref{sec:dynamics}--\ref{sec:sky});
  \item \code{camera.py}: the camera path (Sec.~\ref{sec:camera});
  \item \code{blender\_scene.py}: builds the scene and renders the 3D frames (Sec.~\ref{sec:render});
  \item \code{panel.py}: renders the panel frames and stacks everything into the video.
\end{enumerate}
\code{camera.py}, \code{blender\_scene.py} and \code{panel.py} don't depend on the event: they
read everything from the products directory. The same code also renders GW190521.

\section{Samples}

The cloud is the first $N=200$ of the 500 random draws from the
\code{bilby-SEOBNRv5PHM-reweighted} posterior that are also used for the
posterior-predictive band. The maximum-likelihood sample is added and highlighted.
Each sample keeps its own masses, spins, distance, sky position and 3D orientation.

\section{Dynamics}\label{sec:dynamics}

\paragraph{Model call.} SEOBNRv5PHM (pyseobnr 0.3.5, the version used in the PE run) is
called with exactly the arguments bilby's \code{gwsignal\_binary\_black\_hole} used in the
likelihood: $f_\mathrm{ref}=20$\,Hz, $f_\mathrm{min}=10$\,Hz (lowered by the conditioning),
\code{lmax\_nyquist}$=1$.

\paragraph{Orbit.} The EOB relative orbit in the inertial frame is
\begin{equation}
  \vect{n}(t) = R(t)\,(\cos\phi, \sin\phi, 0)^T ,
\end{equation}
where $\phi$ is the EOB orbital phase and $R(t)$ is the co-precessing$\to$inertial rotation
that pyseobnr uses to twist up the modes. The bodies sit at
\begin{equation}
  \vect{x}_1 = \frac{m_2}{M}\, r\,\vect{n}, \qquad
  \vect{x}_2 = -\frac{m_1}{M}\, r\,\vect{n},
\end{equation}
with $r$ the EOB separation. Lengths are converted to source-frame km with
$G M/c^2/(1+z)$. The spins are pyseobnr's PN-evolved spin vectors, as functions of the EOB
orbital frequency. EOB coordinates depend on the gauge, so positions near merger are
indicative only.

\paragraph{Sizes.} Each black hole is drawn as a sphere with the areal radius of its
Kerr horizon, $\sqrt{A/4\pi} = \sqrt{2 m r_+}$, with $r_+ = m(1+\sqrt{1-\chi^2})$.

\paragraph{Timeline.} Frames are evenly spaced in L1 detector time. The state shown at
$t_\mathrm{det}$ is the one at waveform time
$t_\mathrm{det} - t_\mathrm{geocent} - \Delta t_\mathrm{L1}$, i.e.\ the state that emitted
the signal arriving at L1 at $t_\mathrm{det}$. All samples share this timeline, so each one
merges at its own detector time. Across the cloud, these times span
$[0.21636, 0.21678]$\,s.

\section{Merger and ringdown}\label{sec:ringdown}

\paragraph{Remnant.} From the attachment time $t_\mathrm{attach}$ on, the two bodies are
replaced by the remnant, with pyseobnr's final mass and spin $\chi_f$. The spin points
along $\hat{\vect{J}}$ at attachment, as pyseobnr models it.

\paragraph{Ringing.} In a frame with $z$ along the remnant spin and $\varphi$ measured
from the separation $\vect{n}(t_\mathrm{attach})$, the horizon surface is
\begin{equation}
  r(\theta,\varphi,t) = R_f\left[1 + \operatorname{Re}\!\left(C(t)\,e^{-2i\varphi}\right)
  \sin^2\theta\right], \qquad
  C(t) = 0.3\,\overline{\left(\frac{h(t)}{h(t_\mathrm{attach})}\right)} .
\end{equation}
Here $h$ is a 220\,+\,221 Kerr quasinormal-mode fit to the sample's own $(2,2)$ mode in
the final-$J$ frame, starting at $t_\mathrm{attach}$, with frequencies from the \code{qnm}
package for that sample's $(M_f,\chi_f)$. The $\ell=m=2$ bar therefore starts along the
line joining the two black holes, rotates with the pattern speed $\omega_R/2$, and decays
with the two damping times, using the fitted $221/220$ amplitude ratio ($\approx1.25$) and
relative phase. In Blender this is implemented as two shape keys of the unit sphere,
$x^2-y^2$ and $2xy$, weighted by $\operatorname{Re}C$ and $\operatorname{Im}C$.
The 30\% overall amplitude is \emph{illustrative}. The following are neglected:
spheroidal vs.\ spherical harmonics, retrograde and mirror modes, and any difference
between the horizon's and the waveform's mode content.

\paragraph{Recoil.} The kick velocity $\vect{v}_k$ comes from the NRSur7dq4Remnant fit
(\code{surfinBH}). Its inputs are the spins and orbit at $-100\,M$ from the peak (which is
$t_\mathrm{attach}$ in v5), expressed in the fit's coorbital frame. From attachment on, the
remnant moves as $\vect{x}_f(t) = \vect{v}_k\,(t - t_\mathrm{attach})$. Over the
ringdown window this is a displacement of $\sim\!30$\,km (median) to $\sim\!200$\,km (largest).

\section{Sky frame and orientation}\label{sec:sky}

\paragraph{Frame.} The sky frame has $X$ pointing west, $Y$ north and $Z$ towards Earth. The LAL
source frame is rotated so that the model's polarisation basis maps onto bilby's $(m,n)$
basis for the sample's $\psi$, and the direction to the observer maps onto $+Z$.

\paragraph{$\psi$ vs.\ $\psi+\pi$.} These give identical detector responses: rotating
the whole configuration by $\pi$ about the line of sight is equally allowed. bilby samples
$\psi\in[0,\pi)$. For each sample we show the representative whose sky-projected
$\vect{J}(t\to-\infty)$ lies on the same side as the axial mean of the set. For this posterior
no sample needs to be rotated.

\paragraph{Grid.} The reference grid is the plane through the centre of mass perpendicular
to $\vect{J}(t\to-\infty)$. That direction is obtained by integrating pyseobnr's PN
precession equations backwards from $f_\mathrm{ref}$ (it differs from $\vect{J}(f_\mathrm{ref})$
by $\sim0.1^\circ$ for the max-L sample). For the cloud, we average the samples'
$\vect{J}(-\infty)$ after reflecting through the sky plane those that point away from
us. $\theta_{JN}$ is bimodal ($\approx45^\circ$ or $\approx135^\circ$), and the two families
are mirror images through the sky plane. Along the line of sight they look the same
(same projected ellipses, opposite senses of rotation). The grid only adds light, so it
never hides bodies on its far side.

\section{Camera}\label{sec:camera}

\code{camera.py} blends three shots with smoothstep weights. The target point is
interpolated linearly, the viewing direction by slerp, and the view width and camera
distance logarithmically:
\begin{enumerate}
  \item \emph{From Earth}: on the line of sight, 40\,000\,km away (nearly orthographic),
        framing the orbits with a 10\% margin and zooming in slowly as they shrink;
  \item \emph{Close-up} (in over 39.7--33.7\,s, out over 15.7--8.7\,s before merger in
        video time): follows a cluster of bodies from above its orbital plane, co-rotating
        with it. The cluster is formed by the bodies of the majority $\theta_{JN}$ family
        that are close to the same body of a reference sample;
  \item \emph{Back to Earth}: on the line of sight for the merger and ringdown, framing the
        horizons plus the spin arrows.
\end{enumerate}
When the camera is off the line of sight, a movie-camera icon on that line, joined to the
centre of mass and labelled ``to Earth'', marks the observer's direction. The panel then
shows the off-axis angle.

\section{Rendering}\label{sec:render}

\code{blender\_scene.py} (Blender 5.2, EEVEE, 16 samples per pixel, 1 Blender unit
$=100$\,km) builds one set of objects per sample, keyframed on every frame with constant
interpolation, and uses the \code{solid\_highlight} style:

\begin{itemize}
  \item \textbf{Horizons}: opaque black spheres with coloured Fresnel rims. The heavier
        body is orange and the lighter one blue, and the remnant is green. Since
        $q\approx1$, which body is the heavier one varies across samples.
  \item \textbf{Spin arrows}: along $\vect{\chi}_i$, starting from the centre and
        protruding from the horizon by $400\,\mathrm{km}\times|\chi|$. Each remnant
        arrow has a small flared collar where it leaves the horizon, at the spin pole.
        The $\sin^2\theta$ deformation vanishes there, so the collar stays on the surface
        while the remnant rings. Without the collar, an arrow leaving the face of the
        horizon that faces the camera seems to float in front of it. This is the usual
        case for the remnant, whose spin is only $\sim40^\circ$ from the line of sight.
  \item \textbf{Trails}: thin tubes along each body's past orbit, sampled at 96 points
        per orbit. They are alpha-blended so that they saturate to the body colour,
        independently of the drawing order. Their brightness falls off exponentially
        with the orbital phase elapsed since the body passed there (e-folding 0.3
        orbits), with a smooth cutoff between 0.4 and 0.6 orbits. After the attachment,
        the phase keeps running at the attachment orbital frequency, so the trails fade
        out.
  \item \textbf{Max-likelihood sample}: lighter colours, thicker arrows and longer opaque
        trails.
\end{itemize}

\section{Strain panel}\label{sec:panel}

\code{ppd.py} reproduces the whitened-strain reconstruction of the GW250114 paper
(\code{waveform\_reconstruction.py}): same data dump, PSDs, whitening and
$\sqrt{\beta}$ factor ($\alpha=1/4$, $\beta = 1-5\alpha/8$). The waveforms come from 500
SEOBNRv5PHM samples, generated with the PE settings and per-sample calibration, and the
panel shows their 90\% band. The NRSur7dq4 run's data dump is used because its PSDs are
identical to the SEOBNRv5PHM run's. \code{panel.py} draws the L1 data, the band, the
attachment times and a cursor, in a window that scrolls 13\,s behind and 6\,s ahead of the
cursor (in video time). It then stacks the panel under the 3D frames with ffmpeg.

\section{Checks}\label{sec:checks}

\code{validate\_dynamics.py} checks the following:
\begin{enumerate}
  \item $R(t)\hat z$ coincides with $\hat{\vect{L}}_N(t)$ to $0.002^\circ$. At
        $t_\mathrm{ref}$, $\vect{n}=\hat x$ and $\hat{\vect{L}}_N = \hat z$ (LAL convention).
  \item $h_+, h_\times$ computed from the point-mass orbit with the Cartesian multipole
        formula (mass quadrupole, octupole, current quadrupole) match the model's to
        $\approx20\%$. The residual comes from PN corrections that the leading-order
        formula doesn't include. Other polarisation bases give $>130\%$, and swapping the
        bodies is disfavoured in an unequal-mass ($q=3$) test.
  \item bilby's frequency-domain L1/H1 response matches the time-domain polarisations
        placed at $t_\mathrm{geocent}+\Delta t$ with zero lag.
  \item With NRSur7dq4, \code{ppd.py} reproduces the paper's reconstruction.
  \item Ringdown: the 220\,+\,221 fit of $h_{22}$ in the final-$J$ frame has a relative
        residual of 0.18--0.21 (0.54 with 220 alone). Almost all of it is within the first
        few $M$, where the model's amplitude has a flat peak that no QNM sum reproduces.
        $|A_{221}/A_{220}|\approx1.25$ is stable against the fit length. Running pyseobnr
        with mode output leaves the dynamics bit-identical.
\end{enumerate}

\section{Caveats}

\begin{itemize}
  \item Positions near merger are in EOB coordinates, which depend on the gauge.
  \item The ringdown deformation amplitude is illustrative. Only its time dependence
        (rotation, decay, 221/220 content) comes from the model.
  \item The two bodies switch to the remnant instantaneously at $t_\mathrm{attach}$; no
        common horizon is modelled.
  \item The 3D view uses perspective projection. Only the ``from Earth'' shots, being
        nearly orthographic, correspond to what a distant observer would ``see''.
\end{itemize}

\section*{Reproducing}

\begin{verbatim}
uv sync --python 3.10
.venv/bin/python validate_dynamics.py
.venv/bin/python ppd.py --ifo L1 --include-maxl
.venv/bin/python export_scene.py --n-samples 200
.venv/bin/python camera.py --name cloud200
blender -b --factory-startup -P blender_scene.py -- \
    --manifest products/cloud200.json --outdir renders/cloud200 \
    --style solid_highlight
.venv/bin/python panel.py --name cloud200
\end{verbatim}
\code{blender\_scene.py --frames a:b} (or a comma-separated list) re-renders a subset of the
frames. For example, the ringdown frames were re-rendered with \code{--frames 1520:} after
adding the collars to the remnant spin arrows.

\end{document}
